\chapter*{Preface}
Tome~2 is about structure: groups, rings, fields and their representations. It builds abstract algebra from
symmetry and climbs to two summits. The first is Galois' answer to an old question: there is a formula for the
roots of every equation of degree $4$, but none for the equation $x^5-4x+2=0$. The second is the character table
of a finite group, a small square of numbers that holds most of what there is to know about the group.

It assumes Tomes~0 and~1 (references such as ``Tome~1, Theorem~5.3'' point there). Exercises are graded
\textbf{$\star$} warm-up, \textbf{$\star\star$} for a good afternoon, \textbf{$\star\star\star$} gems, and every exercise has a
hint and a complete solution at the end of its chapter. Each chapter has a companion Python notebook, where the
groups of this book are built, counted and drawn by the computer.

The rules of the earlier tomes hold here too. Every statement is proved, or, if quoted only for culture, labelled
as a remark that says where it is proved. Every solution starts with \emph{How to find it} before the
\emph{Write-up}. ``Pause and think'' boxes ask you to guess before you read on: the guess is the point, not
whether it is right.

The chapters follow two tracks that meet at Galois theory. Chapters~1--5 are about groups, Chapters~6--10
about rings and fields, Chapter~11 about representations, and Chapter~12 about rotations of space.
